\documentclass[12pt,a4paper]{article}

\usepackage{fontawesome5}
\providecommand{\faPencilAlt}{\faIcon{pencil-alt}}

\usepackage{styles/couleurs}
\usepackage{styles/mise_en_page}
\usepackage{styles/boites}
\usepackage{styles/maths}
\usepackage{styles/figures}

\usepackage{enumitem}
\usepackage{array}
\usepackage{tabularx}
\usepackage{amsmath}
\usepackage{amssymb}
\usepackage{graphicx}
\usepackage{tikz}
\usepackage{xparse}

\usetikzlibrary{calc,arrows.meta}

\setlist[itemize]{leftmargin=1.4em,itemsep=0.15em,topsep=0.2em}
\setlist[enumerate]{leftmargin=*,itemsep=0.22em,topsep=0.25em}

\newcommand{\lignepointillee}{%
  \noindent\dotfill\par
  \vspace{0.45em}
}

\newcommand{\lignesreponse}[1]{%
  \begin{tcolorbox}[
    enhanced,
    colback=white,
    colframe=NeutreC!40,
    boxrule=0.5pt,
    arc=1.5mm,
    left=3mm,
    right=3mm,
    top=2mm,
    bottom=1mm,
    before skip=2pt,
    after skip=2pt
  ]
    \foreach \i in {1,...,#1}{\lignepointillee}
  \end{tcolorbox}
}

\newcommand{\imageactivite}[3]{%
  \IfFileExists{#3}{%
    \includegraphics[width=#1,height=#2,keepaspectratio]{#3}%
  }{%
    \fbox{%
      \parbox[c][#2][c]{#1}{%
        \centering\sffamily\small
        Image à intégrer\par
        \vspace{0.3em}
        \texttt{\detokenize{#3}}
      }%
    }%
  }%
}

\newcommand{\titreseq}[1]{%
\begin{center}
  {\sffamily\LARGE\bfseries\color{TitreC}
  Séquence 5 -- Statistiques et proportionnalité}\\[0.2em]
  {\sffamily\large\bfseries\color{SavoirC}#1}
\end{center}
}

\begin{document}

\titreseq{Activité 2 -- Agrandir ou réduire un puzzle}

\remarque{Travail par groupes de trois ou quatre. Chaque groupe traite un niveau.
Le puzzle final est construit en vraie grandeur.}

\begin{center}
\imageactivite
  {0.76\textwidth}
  {0.38\textheight}
  {images/puzzle_brousseau.png}
\end{center}

\section*{Consigne commune}

Toutes les longueurs du puzzle doivent être multipliées par le même nombre.

\begin{itemize}
  \item \textbf{Niveau 1 :} le segment de \(4\) cm doit mesurer \(2\) cm.
  \item \textbf{Niveau 2 :} le segment de \(3\) cm doit mesurer \(4{,}5\) cm.
  \item \textbf{Niveau 3 :} le segment de \(5\) cm doit mesurer \(7\) cm.
\end{itemize}

\section*{Partie A -- Trouver le coefficient}

Pour ton niveau, complète :

\[
\text{nouvelle longueur}
=
\text{longueur du modèle}\times\ldots
\]

Le nombre utilisé est appelé \textbf{coefficient d'agrandissement ou de réduction}.

\lignesreponse{3}

\section*{Partie B -- Préparer la construction}

Complète la ligne correspondant à ton niveau.

\renewcommand{\arraystretch}{1.55}
\begin{tabularx}{\textwidth}{|>{\centering\arraybackslash}p{2.2cm}|*{7}{>{\centering\arraybackslash}X|}}
\hline
\textbf{Longueur modèle} & \(2\) & \(3\) & \(4\) & \(5\) & \(6\) & \(8\) & \(10\)\\
\hline
\textbf{Niveau 1} & & & & & & & \\
\hline
\textbf{Niveau 2} & & & & & & & \\
\hline
\textbf{Niveau 3} & & & & & & & \\
\hline
\end{tabularx}

\section*{Partie C -- Construire et contrôler}

Construis les six pièces, découpe-les puis vérifie qu'elles reconstituent
le même puzzle.

Explique pourquoi il serait faux d'ajouter le même nombre à toutes les longueurs.

\lignesreponse{5}

\end{document}
